\documentclass[11pt,a4paper]{article}

\usepackage[margin=1in]{geometry}
\usepackage{amsmath,amssymb,amsthm,mathtools}
\usepackage{microtype}
\emergencystretch=3em
\usepackage[T1]{fontenc}
\usepackage[utf8]{inputenc}
\usepackage{lmodern}
\usepackage{booktabs}
\usepackage{enumitem}
\usepackage[dvipsnames]{xcolor}
\usepackage{hyperref}
\usepackage[nameinlink,capitalize]{cleveref}

\hypersetup{
  colorlinks=true,
  linkcolor=blue!50!black,
  citecolor=blue!50!black,
  urlcolor=blue!50!black,
  pdftitle={Analytic inversion of two-dimensional Fourier optics in LEEM},
  pdfauthor={Wilson Lau}
}

\newtheorem{theorem}{Theorem}[section]
\newtheorem{lemma}[theorem]{Lemma}
\newtheorem{proposition}[theorem]{Proposition}
\newtheorem{corollary}[theorem]{Corollary}
\theoremstyle{definition}
\newtheorem{definition}[theorem]{Definition}
\theoremstyle{remark}
\newtheorem{remark}[theorem]{Remark}

\newcommand{\R}{\mathbb{R}}
\newcommand{\C}{\mathbb{C}}
\newcommand{\Z}{\mathbb{Z}}
\newcommand{\dd}{\mathrm{d}}
\newcommand{\ee}{\mathrm{e}}
\newcommand{\ii}{\mathrm{i}}
\newcommand{\conj}[1]{\overline{#1}}
\newcommand{\abs}[1]{\lvert#1\rvert}
\newcommand{\norm}[1]{\lVert#1\rVert}
\newcommand{\qap}{q_{\mathrm{ap}}}
\newcommand{\Dz}{\Delta z}
\newcommand{\RFO}{R_{\mathrm{FO}}}
\newcommand{\RFOII}{R_{\mathrm{FO2}}}
\newcommand{\RCTF}{R_{\mathrm{CTF}}}

\title{Analytic inversion of two-dimensional Fourier optics\\
in aberration-corrected LEEM}
\author{Wilson Lau}
\date{26 August 2026}

\begin{document}
\maketitle

\begin{abstract}
We give the closed-form inverse of discrete bilinear Fourier optics
(FO) for a two-dimensional through-focal LEEM series, against the
kernel of Yu, Lau and Altman~\cite{Yu2019}.
Intensity at difference frequency $\xi$ is linear in a lag diagonal of
the rank-1 Gram $X(q,q')=\Psi(q)\conj{\Psi(q')}$.
A vacuum-gauged spectrum is the DC-column factor of $X$.
Subtracting the off-axis remainder reduces each lag to a vacuum
$2\times 2$ that Cramer's rule inverts at vanishing Tikhonov parameter
whenever the Gram determinant is nonzero.
The remainder vanishes on one-, two- and three-mode objects and on
non-collinear three-wave objects $\{0,u,v\}$.
Under perfect coherence and pure defocus the kernel is a cisoid whose
frequency depends on $q$ only through $q\cdot\xi$; equally spaced
defoci are a Vandermonde system for the iso-$\omega$ fiber masses, which
become Gram entries on transversal support.
Householder partners across $\xi$ share a node, so they remain one
mass.
The identities are machine-checked in Lean~4.
This is not an inverse of unrestricted Yu $\RFOII$ on a general
lattice: the DC slice of a unimodular diagonal kernel has rank one,
and radial pure defocus cannot separate perpendicular frequencies.
\end{abstract}

\tableofcontents

\section{Purpose and standing hypotheses}
\label{sec:purpose}

The forward map is the FO kernel of Yu, Lau and
Altman~\cite{Yu2019}, in the line of Pang~\emph{et
al.}~\cite{Pang2009} and Yu~\emph{et al.}~\cite{Yu2017}.
The companion inverse note~\cite{Lau2026inverse} recommends a
\emph{linearized} estimator: multi-defocus Tikhonov on the CTF slice
$\RFO(\mathbf{q},0,\Dz)$, optionally followed by Gauss--Newton on the
bilinear residual.
That map is the Fr\'echet derivative of FO at a strong specular beam.
The present note inverts the \emph{bilinear} map itself, in closed
form, on a finite abelian frequency group.
It is a derivation of identities, not an experimental protocol and not
a substitute for~\cite{Lau2026inverse} on a general micrograph.

Standing discrete data.
Let $G$ be a finite abelian group of sampled spatial frequencies
(in Lean: any \texttt{AddCommGroup} with \texttt{Fintype}) and let
$\kappa$ index a finite through-focal series.
The object spectrum is $\Psi:G\to\C$.
A kernel at one defocus is a map $R:G\times G\to\C$; a stack is
written $R_k$.
The Yu 2D kernel $\RFOII$ is sampled on $G$ by an embedding
$\iota:G\to\R^2$ (\texttt{sampledR}).
Wave aberration $\chi$ is in waves, $W=\ee^{2\pi\ii\chi}$, with
$\mathbf{q}=\boldsymbol{\alpha}/\lambda$ after acceleration~\cite{Tromp2012,Schramm2012,Lau2026proofs}.

The reconstruction proceeds in seven steps
(\cref{sec:gram,sec:gauge,sec:split,sec:small,sec:kernel2,sec:fiber,sec:recipe}).
Each step is an identity; Lean names are given after the statement.
No probability space, sampling theorem, or Gauss--Newton convergence
estimate is used.

\section{Step 1: lift intensity to a Gram contraction}
\label{sec:gram}

The discrete Fourier coefficient of recorded intensity at difference
frequency $\xi\in G$ is
\begin{equation}
\label{eq:ihat}
\hat I(\xi)
=\sum_{q\in G}
\Psi(q)\,R(q,q-\xi)\,\conj{\Psi(q-\xi)}.
\end{equation}
Define the rank-1 Gram
\begin{equation}
\label{eq:gram}
X(q,q')
:=\Psi(q)\,\conj{\Psi(q')}.
\end{equation}
Reassociating~\eqref{eq:ihat} yields the linear contraction
\begin{equation}
\label{eq:ihatgram}
\hat I(\xi)
=\sum_{q\in G}
R(q,q-\xi)\,X(q,q-\xi).
\end{equation}
Thus each fixed-$\xi$ through-focal series is a linear image of the
lag diagonal $q\mapsto X(q,q-\xi)$.
Inverting bilinear FO splits into (i)~inverting those linear slices
for $X$ and (ii)~factoring the Gram for $\Psi$.

Lean: \texttt{ihat}, \texttt{gram}, \texttt{ihat\_eq\_gram}.

\begin{lemma}[Hermitian intensity]
\label{lem:herm}
If $\conj{R(q,q')}=R(q',q)$ for all $q,q'$, then
$\conj{\hat I(\xi)}=\hat I(-\xi)$.
\end{lemma}
Lean: \texttt{ihat\_hermitian}, and \texttt{R\_FO2\_hermitian} for the
Yu kernel (nonnegative energy spread).

\begin{lemma}[Rank-1 pupil]
\label{lem:rank1}
If $R(q,q')=P(q)\conj{P(q')}$, then $\hat I$ is the Gram
autocorrelation of $q\mapsto\Psi(q)P(q)$.
\end{lemma}
Lean: \texttt{ihat\_of\_rank1}.
Perfect coherence of $\RFOII$ is of this form
(\texttt{R\_FO2\_eq\_rank1\_of\_perfect},
\texttt{ihat\_sampledR\_perfect}).

\section{Step 2: quotient the global phase}
\label{sec:gauge}

Intensity is invariant under $\Psi\mapsto\ee^{\ii\theta}\Psi$:
\begin{equation}
\label{eq:gauge}
X(\ee^{\ii\theta}\Psi)=X(\Psi),
\qquad
\hat I(R,\ee^{\ii\theta}\Psi)=\hat I(R,\Psi).
\end{equation}
Lean: \texttt{gram\_eq\_of\_phase}, \texttt{ihat\_gauge}.

The inverse therefore recovers $\Psi$ only in a gauge.
The vacuum gauge takes the specular coefficient real and nonnegative,
\begin{equation}
\label{eq:vacuum}
\hat\Psi(q)
=\frac{X(q,0)}{\sqrt{\norm{X(0,0)}}},
\end{equation}
which uses only the DC column of $X$.

\begin{theorem}[Vacuum-gauge recovery]
\label{thm:vacuum}
If $\Psi(0)\neq 0$, $\operatorname{Im}\Psi(0)=0$ and
$\operatorname{Re}\Psi(0)\ge 0$, then~\eqref{eq:vacuum} applied to
$X(\Psi)$ returns $\Psi$.
\end{theorem}
Lean: \texttt{vacuumGaugePsi}, \texttt{vacuumGaugePsi\_recovers},
\texttt{analyticPsi\_recovers}.
Without the phase convention one obtains only a unimodular factor
$\conj{\Psi(0)}/\abs{\Psi(0)}$
(\texttt{vacuumGaugePsi\_gram}).

\begin{corollary}[Uniqueness from the DC column]
\label{cor:dccol}
If two vacuum-gauged spectra (each with nonzero real-nonnegative
specular coefficient) have the same DC column
$q\mapsto X(q,0)$, they coincide.
\end{corollary}
Lean: \texttt{analyticPsi\_unique\_of\_dc\_column}.
A full Gram is stronger than needed
(\texttt{analyticPsi\_unique}).

Consequently it is enough to recover $X(\cdot,0)$, including the
scalar $X(0,0)$, and then apply~\eqref{eq:vacuum}.

\section{Step 3: split each lag into vacuum $2\times 2$ plus remainder}
\label{sec:split}

For $\xi\neq 0$ the sum~\eqref{eq:ihatgram} isolates the two vacuum
pairs $(\xi,0)$ and $(0,-\xi)$:
\begin{equation}
\label{eq:split}
\hat I(\xi)
=R(\xi,0)\,X(\xi,0)
+R(0,-\xi)\,X(0,-\xi)
+\mathcal{B}(R,\Psi;\xi),
\end{equation}
where the bilinear remainder is the leftover sum
\begin{equation}
\label{eq:B}
\mathcal{B}(R,\Psi;\xi)
=\sum_{q\notin\{0,\xi\}}
R(q,q-\xi)\,X(q,q-\xi).
\end{equation}
Lean: \texttt{ihat\_dc\_split}, \texttt{bilinearRemainder},
\texttt{ihat\_sub\_remainder}.

Subtracting a known remainder produces the vacuum model
\begin{equation}
\label{eq:twocol}
\hat I(\xi)-\mathcal{B}(R,\Psi;\xi)
=h\,U+g\,V,
\end{equation}
with through-focal columns $h_k=R_k(\xi,0)$,
$g_k=R_k(0,-\xi)$ and unknowns
$U=X(\xi,0)$, $V=X(0,-\xi)$.

Write
\begin{equation}
\label{eq:gramdet}
\Delta
=\Bigl(\sum_k\abs{h_k}^2\Bigr)
 \Bigl(\sum_k\abs{g_k}^2\Bigr)
-\Bigl\lvert\sum_k\conj{h_k}\,g_k\Bigr\rvert^2
\end{equation}
for the Gram determinant of the $2\times 2$ at vanishing Tikhonov
parameter $\alpha=0$.

\begin{theorem}[Cramer recovery of the vacuum pair]
\label{thm:cramer}
If $\xi\neq 0$ and $\Delta\neq 0$, then the $\alpha=0$ Cramer formula
applied to the remainder-corrected stack recovers
$\bigl(X(\xi,0),X(0,-\xi)\bigr)$ exactly.
Equal remainder-corrected stacks imply equal vacuum pairs.
\end{theorem}
Lean: \texttt{twoColumn\_recovers},
\texttt{analyticPair\_of\_remainder},
\texttt{dcPair\_unique\_of\_remainder},
\texttt{twoColumn\_injective}, \texttt{gramDet}.

The linearized stage-1 pair estimator of~\cite{Lau2026inverse} is the
same Cramer formula with $\mathcal{B}$ replaced by zero
(\texttt{stage1Pair\_eq\_analyticPair\_linearized}).
Exact vacuum-column recovery therefore requires either knowledge of
$\mathcal{B}$ or a support that forces $\mathcal{B}=0$
(\cref{sec:small}).

\section{Step 4: invert when the remainder vanishes}
\label{sec:small}

Support restrictions make~\eqref{eq:B} identically zero, so
\cref{thm:cramer} applies to the raw intensities.

\subsection{One Fourier mode}
\label{sec:onemode}

If $\Psi(q)=0$ for all $q\neq 0$, then
\begin{equation}
\hat I(\xi)
=
\begin{cases}
R(0,0)\,X(0,0), & \xi=0,\\
0, & \xi\neq 0.
\end{cases}
\end{equation}
Lean: \texttt{ihat\_of\_dc\_support}.

\subsection{Two Fourier modes}
\label{sec:twomode}

Let $\operatorname{supp}\Psi\subseteq\{0,\xi\}$ with $\xi\neq 0$ and
$2\xi\neq 0$ (so $-\xi$ is not a mode).
Then $\mathcal{B}=0$, the conjugate branch $X(0,-\xi)$ vanishes, and
\begin{equation}
\label{eq:twomode}
\hat I(\xi)=R(\xi,0)\,X(\xi,0).
\end{equation}
If $\sum_k\abs{R_k(\xi,0)}^2\neq 0$, scalar Tikhonov at $\alpha=0$
recovers $X(\xi,0)$.
Lean: \texttt{bilinearRemainder\_twoMode}, \texttt{ihat\_twoMode},
\texttt{analyticPair\_twoMode}.

The two amplitudes satisfy the quadratic
$s(T-s)=\abs{Z}^2$ with
$T=\abs{\Psi(0)}^2+\abs{\Psi(\xi)}^2$ and $Z=X(\xi,0)$.
The roots are
\begin{equation}
\label{eq:quad}
s_\pm
=\frac{T\pm\sqrt{T^2-4\abs{Z}^2}}{2};
\end{equation}
the specular (larger-DC) root is $s_+$ whenever
$\abs{\Psi(\xi)}^2\le\abs{\Psi(0)}^2$.
Lean: \texttt{twoModePlus}, \texttt{twoModeMinus},
\texttt{twoMode\_roots\_pair}, \texttt{twoMode\_specular}.
The DC intensity is
$R(0,0)X(0,0)+R(\xi,\xi)X(\xi,\xi)$
(\texttt{ihat\_twoMode\_dc}).

\subsection{Three collinear modes}
\label{sec:threemode}

Let $\operatorname{supp}\Psi\subseteq\{-\xi,0,\xi\}$ with
$\xi\neq 0$, $2\xi\neq 0$ and $3\xi\neq 0$.
Every leftover index in~\eqref{eq:B} hits a missing mode, so
$\mathcal{B}=0$ and
\begin{equation}
\label{eq:threemode}
\hat I(\xi)
=R(\xi,0)\,X(\xi,0)+R(0,-\xi)\,X(0,-\xi).
\end{equation}
\Cref{thm:cramer} applies to the raw stack.
Lean: \texttt{bilinearRemainder\_threeMode},
\texttt{ihat\_threeMode\_axis}, \texttt{analyticPair\_threeMode}.

\subsection{Two-axis three-wave object}
\label{sec:twoaxis}

Let $\operatorname{supp}\Psi\subseteq\{0,u,v\}$ with
$u\neq 0$, $v\neq 0$, $u\neq v$, $2u\neq 0$, $2u\neq v$ and
$v\neq -u$.
(The group hypotheses are the precise content of ``non-collinear'' on
a general abelian group: $v-u\notin\{0,u,v\}$ and $-\,u\notin\{0,u,v\}$.)
Then at frequency $u$ the remainder and conjugate branch drop,
\begin{equation}
\label{eq:twoaxis}
\hat I(u)=R(u,0)\,X(u,0).
\end{equation}
Scalar Tikhonov at $\alpha=0$ recovers $X(u,0)$ whenever
$\sum_k\abs{R_k(u,0)}^2\neq 0$.
The swapped hypotheses at $v$ recover $X(v,0)$.
{\raggedright
Lean: \texttt{bilinearRemainder\_twoAxis}, \texttt{ihat\_twoAxis},
\texttt{ihat\_twoAxis\_v}, \texttt{analyticPair\_twoAxis},
\texttt{analyticPair\_twoAxis\_v}, \texttt{twoAxis\_dcCol\_unique},
\texttt{twoAxis\_dcCol\_unique\_v}.
\par}
The DC intensity is the sum of the three diagonal Gram entries
(\texttt{ihat\_twoAxis\_dc}).

Together with a determination of $X(0,0)$ (\cref{sec:dcenergy}) this
is the full DC column of a three-wave object, hence $\Psi$ by
\cref{cor:dccol}.

\section{Step 5: pass to the 2D Yu kernel}
\label{sec:kernel2}

The working 2D kernel uses the disk aperture, radial
$\chi_{S2}$, and vector tilt
$\mathbf{a}=\nabla\chi_S(\mathbf{q})-\nabla\chi_S(\mathbf{q}')$, so
the spatial envelope $E_S$ sees $\mathbf{q}\cdot\mathbf{q}'$~\cite{Yu2019,Lau2026proofs}.
On the CTF slice
\begin{equation}
\label{eq:axis}
\RFOII(\mathbf{q},0,\Dz)
=\RFO\bigl(\norm{\mathbf{q}},0,\Dz\bigr);
\end{equation}
off axis $\RFOII$ is not the radial 1D kernel evaluated at
$\bigl(\norm{\mathbf{q}},\norm{\mathbf{q}'}\bigr)$.
Lean: \texttt{R\_FO2}, \texttt{R\_FO2\_eq\_R\_FO\_axis}.

\begin{theorem}[Pure-defocus cisoid]
\label{thm:cisoid}
Under perfect coherence ($q_{\mathrm{ill}}=\Delta E=0$) and
$C_3=C_5=0$,
\begin{equation}
\label{eq:cisoid}
\RFOII(\mathbf{q},\mathbf{q}',\Dz)
=A(\mathbf{q})A(\mathbf{q}')\,
\ee^{\ii\,\omega(\mathbf{q},\mathbf{q}')\,\Dz},
\end{equation}
with aperture indicators $A$ and instantaneous frequency
\begin{equation}
\label{eq:omega}
\omega(\mathbf{q},\mathbf{q}')
=\pi\lambda\bigl(\norm{\mathbf{q}}^2-\norm{\mathbf{q}'}^2\bigr)
=\pi\lambda\bigl(2\langle\mathbf{q},\mathbf{q}-\mathbf{q}'\rangle
-\norm{\mathbf{q}-\mathbf{q}'}^2\bigr).
\end{equation}
\end{theorem}
Lean: \texttt{R\_FO2\_pureDefocus\_perfect}, \texttt{defocusOmega},
\texttt{defocusOmega\_eq\_inner}.

At fixed lag $\boldsymbol{\xi}=\mathbf{q}-\mathbf{q}'$,~\eqref{eq:omega}
depends on $\mathbf{q}$ only through the projection
$\langle\mathbf{q},\boldsymbol{\xi}\rangle$.
Frequencies with the same projection (and matching aperture) are
indistinguishable.

\begin{definition}[Householder reflection]
\label{def:refl}
For $\boldsymbol{\xi}\neq\mathbf{0}$,
\begin{equation}
\label{eq:refl}
\operatorname{refl}_{\boldsymbol{\xi}}(\mathbf{q})
=\frac{2\langle\mathbf{q},\boldsymbol{\xi}\rangle}{\norm{\boldsymbol{\xi}}^2}\,
\boldsymbol{\xi}-\mathbf{q}.
\end{equation}
\end{definition}
The reflection preserves $\langle\cdot,\boldsymbol{\xi}\rangle$, Euclidean
norms, $\norm{\mathbf{a}}$, and $\omega$.
Lean: \texttt{reflectLine}, \texttt{defocusOmega\_reflect},
\texttt{R\_FO2\_pureDefocus\_perp}.

Sampling $\RFOII$ on $G$ via $\iota$ yields
$R^\iota_{\Dz}(q,q')=\RFOII(\iota q,\iota q',\Dz)$.
Bilinear FO on that group is~\eqref{eq:ihat} with this kernel
(\texttt{sampledR}, \texttt{ihat\_sampledR}).

\section{Step 6: invert iso-\texorpdfstring{$\omega$}{omega} fibers by a Vandermonde system}
\label{sec:fiber}

Under the hypotheses of \cref{thm:cisoid}, substitute~\eqref{eq:cisoid}
into~\eqref{eq:ihatgram} at lag $\xi$.
Group $q\in G$ by a fiber index $\operatorname{idx}:G\to\{0,\dots,n-1\}$
constant on level sets of
$\omega\bigl(\iota q,\iota(q-\xi)\bigr)$.
The aperture-weighted fiber masses are
\begin{equation}
\label{eq:Mj}
M_j
=\sum_{\operatorname{idx}(q)=j}
A(\iota q)\,A\bigl(\iota(q-\xi)\bigr)\,X(q,q-\xi).
\end{equation}
The through-focal coefficient is then the cisoid mixture
\begin{equation}
\label{eq:mixture}
\hat I(\xi,\Dz)
=\sum_{j=0}^{n-1}
M_j\,\ee^{\ii\omega_j\Dz}.
\end{equation}
Lean: \texttt{apertureGram}, \texttt{cisoidMass},
\texttt{ihat\_eq\_cisoidMass}.

Take $n$ equally spaced defoci $\Dz_k=\delta k$, $k=0,\dots,n-1$,
and write $z_j=\ee^{\ii\omega_j\delta}$.
Equation~\eqref{eq:mixture} becomes the Vandermonde system
\begin{equation}
\label{eq:vandermonde}
\hat I(\xi,\delta k)
=\sum_{j}M_j\,z_j^k.
\end{equation}
If $j\mapsto z_j$ is injective, the nodes are distinct, the
Vandermonde matrix $V_{ij}=z_i^{j}$ is invertible, and
\begin{equation}
\label{eq:recM}
\mathbf{M}
=\mathbf{y}\,V(z)^{-1}
\end{equation}
recovers the masses from the samples
$y_k=\hat I(\xi,\delta k)$.
{\raggedright
Lean: \texttt{cisoidNode}, \texttt{recoveredMasses},
\texttt{cisoidMasses\_unique}, \texttt{recoveredMasses\_eq},
\texttt{ihat\_eq\_vandermondeRow},
\texttt{recoveredMasses\_eq\_cisoidMass}.
\par}

Householder partners~\eqref{eq:refl} share a node, so they contribute
to a single $M_j$.
The inverse of~\eqref{eq:vandermonde} recovers \emph{fiber masses},
not individual lattice points on an iso-$\omega$ line.

\section{Step 7: transversal support turns masses into Gram entries}
\label{sec:transversal}

\begin{definition}[Transversal support]
\label{def:trans}
The pair $(\Psi,\xi)$ is transversal for $\operatorname{idx}$ if each
fiber carries at most one occupied pair:
whenever $\operatorname{idx}(q)=\operatorname{idx}(q')$ and
$\Psi(q)$, $\Psi(q-\xi)$, $\Psi(q')$, $\Psi(q'-\xi)$ are all nonzero,
one has $q=q'$.
\end{definition}
Lean: \texttt{transversalSupport}.

\begin{theorem}[Mass equals Gram]
\label{thm:trans}
Assume \cref{thm:cisoid}, distinct cisoid nodes, transversal support,
and $A=1$ at every occupied frequency.
If $\Psi(q)\neq 0$ and $\Psi(q-\xi)\neq 0$, then
\begin{equation}
\label{eq:Mgram}
M_{\operatorname{idx}(q)}
=X\bigl(q,q-\xi\bigr).
\end{equation}
In particular, if $\Psi(\xi)\neq 0$ and $\Psi(0)\neq 0$, the vacuum
column entry is the mass of the fiber containing $(\xi,0)$:
\begin{equation}
\label{eq:Mvac}
M_{\operatorname{idx}(\xi)}
=X(\xi,0).
\end{equation}
Equal through-focal stacks then imply equal Gram entries at every
commonly occupied pair.
\end{theorem}
{\raggedright
Lean: \texttt{cisoidMass\_eq\_gram\_of\_transversal},
\texttt{recoveredMasses\_eq\_gram\_of\_transversal},
\texttt{recoveredMasses\_eq\_gram\_vacuum},
\texttt{gram\_eq\_of\_transversal\_ihat}.
\par}

If both Householder partners are occupied, \cref{def:trans} fails on
that fiber: the mass is the sum of two Gram entries and cannot be
split.
That is the geometric content of \cref{thm:cisoid} for the inverse.

\subsection{Recovering \texorpdfstring{$X(0,0)$}{X(0,0)}}
\label{sec:dcenergy}

Off-DC fibers never produce the scalar $X(0,0)$.
If $R(q,q)=1$ for all $q$ and $\hat I(0)\neq 0$, the DC intensity is
the total energy $T=\sum_q\abs{\Psi(q)}^2$.
Write $a=X(0,0)=\abs{\Psi(0)}^2$ and
$S=\sum_{q\neq 0}\abs{X(q,0)}^2$.
Then $S=a(T-a)$, so $a$ is a root of the same quadratic as in
\cref{sec:twomode},
\begin{equation}
\label{eq:recDC}
a(T-a)=S,
\qquad
a_\pm
=\frac{T\pm\sqrt{T^2-4S}}{2}.
\end{equation}
The specular choice is the larger root when the DC beam dominates.
The identity
\begin{equation}
\label{eq:recDCcheck}
X(0,0)
=\frac{\sum_q\norm{X(q,0)}^2}{\hat I(0)}
\end{equation}
holds for the \emph{full} true DC column (including $a$ itself); it
is a consistency check, not an inverse for the missing entry.
Lean: \texttt{recoveredDC}, \texttt{recoveredDC\_gram},
\texttt{twoMode\_specular}.

The DC \emph{slice} $q\mapsto X(q,q)$ of a unimodular diagonal kernel
is not injective: it has rank one.
Lean: \texttt{dcSlice\_not\_injective}.
Off-DC slice injectivity recovers lag diagonals with $\xi\neq 0$
(\texttt{offDiagGram\_eq\_of\_sliceInjective}) but does not invert
$\xi=0$.

\section{Reconstruction recipe}
\label{sec:recipe}

The seven identities assemble into the following procedure.
Input: a through-focal stack on a finite frequency group, an embedding
$\iota$, and (as appropriate) support or fiber-index data.
Output: the vacuum-gauged spectrum $\hat\Psi$ when the DC column is
determined, else the fiber masses (or vacuum pairs) that the data
identify.

\begin{enumerate}[leftmargin=*,itemsep=0.35em]
\item \textbf{Form $\hat I(\xi,\Dz_k)$.}
After alignment, discrete-Fourier transform each plane.
Work with~\eqref{eq:ihat}; modes $\norm{\iota q}>\qap$ are
identically invisible.

\item \textbf{Fix the gauge.}
Seek $\Psi$ with $\Psi(0)$ real and nonnegative
(\cref{thm:vacuum}).
A global phase will never be recovered~\eqref{eq:gauge}.

\item \textbf{Choose the lag inverse} according to what is known.

\begin{enumerate}[label=\alph*.,itemsep=0.25em]
\item \emph{Known remainder or small support.}
If $\mathcal{B}$ is known, or if $\Psi$ is supported on
$\{0\}$, $\{0,\xi\}$, $\{-\xi,0,\xi\}$, or $\{0,u,v\}$ as in
\cref{sec:small}, apply \cref{thm:cramer} or the scalar
Tikhonov identity~\eqref{eq:twomode}--\eqref{eq:twoaxis} at
$\alpha=0$, provided the relevant energy $\Delta$ or
$\sum\abs{h_k}^2$ is nonzero.
This recovers the vacuum column entries that the support sees.

\item \emph{Pure defocus, perfect coherence, unknown general
support.}
Assume $C_3=C_5=0$ and \cref{thm:cisoid}.
Index $q$ by $\omega(\iota q,\iota(q-\xi))$, sample $n$ equally
spaced defoci, and right-multiply by the inverse Vandermonde
matrix~\eqref{eq:recM}.
This recovers fiber masses $M_j$, not a unique $q$ on each
iso-$\omega$ line.

\item \emph{Transversal support inside the aperture.}
If in addition \cref{def:trans} holds and $A=1$ on the occupied
set, each recovered mass is a Gram entry~\eqref{eq:Mgram}.
The vacuum column is~\eqref{eq:Mvac} at every occupied $\xi$.
\end{enumerate}

\item \textbf{Assemble the DC column.}
Collect every recovered off-DC entry $X(\xi,0)$.
If $R(q,q)=1$, determine $X(0,0)$ from the quadratic~\eqref{eq:recDC}
(two-mode: $S=\abs{X(\xi,0)}^2$; many modes: $S$ is the sum of
squared off-DC column entries).
Equation~\eqref{eq:recDCcheck} then checks consistency of the
completed column.

\item \textbf{Factor the Gram.}
Apply~\eqref{eq:vacuum}.
By \cref{cor:dccol} the vacuum-gauged object is unique among
spectra sharing that DC column.
\end{enumerate}

The linearized Schiske/Wiener map of~\cite{Lau2026inverse} is step~3a
with $\mathcal{B}:=0$ and without a support hypothesis: it is
first-order consistent at vacuum and biased for strong phase.

\section{What this does not invert}
\label{sec:not}

The analytic inverse is closed-form under remainder knowledge,
vanishing remainder (small-mode or two-axis support), or transversal
support with distinct cisoid nodes.
It is \emph{not} an inverse of unrestricted 2D Yu $\RFOII$ on a
general lattice, for two independent reasons already present in the
kernel.

\begin{enumerate}[leftmargin=*]
\item \textbf{DC rank one.}
A unimodular diagonal kernel has a rank-1 DC slice
(\texttt{dcSlice\_not\_injective}).
Matching off-DC lag diagonals do not determine $X$ on the
diagonal $\xi=0$.

\item \textbf{Projection degeneracy.}
Pure-defocus $\omega$ depends on $q$ only through
$q\cdot\xi$~\eqref{eq:omega}.
Householder partners~\eqref{eq:refl} share a cisoid node and
remain a single fiber mass.
No choice of equally spaced $\Dz_k$ splits them.
\end{enumerate}

Statistical noise models, continuum Fourier transforms versus the
array DFT, wall-clock FFT complexity, and local convergence of
Gauss--Newton remain informal, as in~\cite{Lau2026inverse}.

\section{Machine-checked fragment}
\label{sec:lean}

The identities of
\cref{sec:gram,sec:gauge,sec:split,sec:small,sec:kernel2,sec:fiber,sec:transversal}
are encoded in Lean~4~\cite{deMoura2021,mathlib2020} against the FO
kernel of~\cite{Yu2019,Lau2026proofs}.
The statements are finite-dimensional.
Sound proofs depend only on
\texttt{propext}, \texttt{Classical.choice} and \texttt{Quot.sound};
there is no \texttt{sorry}.
Full identifier lists are in the theorem file
\texttt{docs/LINEAR\_INVERSE.md} (T10--T15).

\begin{center}
\small
\begin{tabular}{@{}ll@{}}
\toprule
Step & Lean modules \\
\midrule
1--3 & \texttt{LeemFO/Inverse/Gram.lean}, \texttt{Analytic.lean} \\
4 & \texttt{LeemFO/Inverse/SmallMode.lean} \\
5 & \texttt{LeemFO/Forward/Kernel2.lean}, \texttt{Inverse/Degeneracy.lean} \\
6--7 & \texttt{LeemFO/Inverse/Fiber.lean} \\
\bottomrule
\end{tabular}
\end{center}

\section{Conclusions}
\label{sec:end}

Discrete bilinear FO is a linear map on lag diagonals of a rank-1
Gram, followed by a vacuum-gauge factorization.
The vacuum $2\times 2$ inverts in closed form once the off-axis
remainder is known or killed by support; a pure-defocus through-focal
series inverts iso-$\omega$ fiber masses by a Vandermonde system; and
transversal support promotes those masses to Gram entries.
The same kernel that prevents an unrestricted lattice inverse---a
rank-one DC slice and Householder-glued cisoid nodes---is what the
identities use.
The reconstruction of~\cref{sec:recipe} is therefore the analytic 2D
FO inverse, checked against Yu~\emph{et al.}~\cite{Yu2019}, not a
claim that every 2D lattice object is identifiable from a radial
defocus series.

\begin{thebibliography}{99}

\bibitem{Yu2019}
K.~M. Yu, K.~L.~W. Lau, M.~S. Altman,
\emph{Fourier optics of image formation in aberration-corrected LEEM},
Ultramicroscopy \textbf{200} (2019) 160--168.

\bibitem{Yu2017}
K.~M. Yu, A.~Locatelli, M.~S. Altman,
\emph{Comparing Fourier optics and contrast transfer function modeling
of image formation in low energy electron microscopy},
Ultramicroscopy \textbf{183} (2017) 109--116.

\bibitem{Pang2009}
A.~B. Pang, Th.~M\"uller, M.~S. Altman, E.~Bauer,
\emph{Fourier optics of image formation in LEEM},
J.\ Phys.: Condens.\ Matter \textbf{21} (2009) 314006.

\bibitem{Schramm2012}
S.~M. Schramm, A.~B. Pang, M.~S. Altman, R.~M. Tromp,
\emph{A contrast transfer function approach for image calculations
in standard and aberration-corrected LEEM and PEEM},
Ultramicroscopy \textbf{115} (2012) 88--108.

\bibitem{Tromp2012}
R.~M. Tromp, S.~M. Schramm,
\emph{Optimization and stability of the contrast transfer function
in aberration-corrected electron microscopy},
Ultramicroscopy \textbf{125} (2013) 72--80.

\bibitem{Lau2026proofs}
W.~Lau,
\emph{Closed-form envelopes of the Fourier-optics kernel in
aberration-corrected LEEM},
unpublished (2026).

\bibitem{Lau2026inverse}
W.~Lau,
\emph{Through-focal inversion of LEEM intensities with the
Fourier-optics kernel},
unpublished (2026).

\bibitem{deMoura2021}
L.~de~Moura, S.~Ullrich,
\emph{The Lean 4 theorem prover and programming language},
in \emph{Automated Deduction -- CADE 28},
LNCS \textbf{12699} (Springer, Cham, 2021), pp.~625--635.

\bibitem{mathlib2020}
The mathlib Community,
\emph{The Lean mathematical library},
in \emph{Proc.\ 9th ACM SIGPLAN International Conference on
Certified Programs and Proofs} (ACM, New York, 2020), pp.~367--381.

\end{thebibliography}

\end{document}
